% Rate 1; chosen arrival times 0.6, 1.8, 3.2 (illustration, not simulation).
\begin{tikzpicture}[x=21mm,y=9mm,
  every node/.style={font=\figurefont\footnotesize,text=FigureInk}]
  \draw[->,FigureStroke,line width=.6pt] (0,-1.55)--(0,3.7);
  \draw[->,FigureStroke,line width=.6pt] (0,0)--(4.35,0) node[right] {\(t\)};
  \foreach \t in {1,2,3,4} {
    \draw[FigureMuted] (\t,-.06)--(\t,.06);
    \node[below] at (\t,-.1) {\(\t\)};
  }
  \foreach \y in {-1,1,2,3} {
    \draw[FigureMuted] (-.04,\y)--(.04,\y);
    \node[left] at (-.05,\y) {\(\y\)};
  }
  \draw[FigureBlue,line width=1pt] (0,0)--(.6,0)--(.6,1)
    --(1.8,1)--(1.8,2)--(3.2,2)--(3.2,3)--(4,3);
  \draw[FigureRed,dashed,line width=1pt] (0,0)--(.6,-.6)--(.6,.4)
    --(1.8,-.8)--(1.8,.2)--(3.2,-1.2)--(3.2,-.2)--(4,-1);
  \foreach \t/\n in {.6/1,1.8/2,3.2/3}
    \fill[FigureBlue] (\t,\n) circle (1.5pt);
  \node[anchor=west,text=FigureBlue] at (3.3,3.35) {\(N_t\)};
  \node[anchor=west,text=FigureRed] at (3.1,-1.65) {\(N_t-t\)};
  \node[align=center] at (2,-2.6)
    {跳跃量同为\(1\)；补偿只加入斜率\(-1\)的连续漂移};
\end{tikzpicture}
